\chapter{Caja de herramientas}
\label{cap:herramientas}

Este capítulo repasa las cinco ideas matemáticas que el método usa de forma
intensiva. Si alguna ya resulta familiar, conviene saltarla; si todas resultan
familiares, conviene saltar el capítulo entero y volver solo cuando el texto
remita aquí.

\section{Derivadas de funciones entre espacios de secuencias}

Una \emph{secuencia} no es más que una lista de números indexada por el tiempo:
$\mathbf{x} = (x_0, x_1, x_2, \dots)$. En la práctica siempre la truncamos en un
horizonte $T$, de modo que $\mathbf{x} \in \R^T$ es un vector columna.

Si una función $h$ toma una secuencia y devuelve otra secuencia,
$\mathbf{y} = h(\mathbf{x})$ con $\mathbf{y}\in\R^T$, entonces su derivada en un
punto es una matriz $T\times T$:
\begin{equation}
\J_{t,s} = \frac{\partial y_t}{\partial x_s}.
\end{equation}
Eso es todo lo que significa ``jacobiano secuencial''. No hay nada más profundo
en el nombre. Lo único inusual es que estamos derivando respecto de una variable
\emph{con fecha}: $\partial y_t/\partial x_s$ pregunta cuánto cambia el producto
en la fecha $t$ si el insumo cambia en la fecha $s$, manteniendo fijo el insumo
en todas las demás fechas.

\begin{idea}
Leer un jacobiano secuencial es leer dos cosas distintas según el eje.
Una \textbf{columna} $s$ fija es una función impulso-respuesta: la trayectoria
completa de $y$ ante un cambio en $x$ en la fecha $s$. Una \textbf{fila} $t$
fija es una función de ponderación: cuánto pesa cada fecha del insumo en el
valor del producto en $t$.
\end{idea}

Dos propiedades de estas matrices aparecerán una y otra vez.

\paragraph{Causalidad y anticipación} Si $y_t$ no puede depender de lo que pasa
después de $t$, entonces $\J_{t,s} = 0$ para $s > t$ y la matriz es triangular
inferior. Eso ocurre en bloques donde no hay expectativas. Pero cuando los
agentes son \emph{prospectivos} --- ahorran hoy porque anticipan una tasa más
alta mañana --- el triángulo superior es distinto de cero, y de hecho es donde
vive casi toda la economía interesante del modelo.

\paragraph{Estructura de Toeplitz} Una matriz es de Toeplitz si solo depende de
$t-s$: $\J_{t,s} = g(t-s)$, es decir, es constante a lo largo de las diagonales.
Esto equivale a decir que la respuesta a una noticia solo depende de cuán lejana
sea la noticia, no de en qué fecha del calendario estemos. Los jacobianos de los
modelos heterogéneos \emph{no} son exactamente de Toeplitz --- justamente porque
la distribución cambia --- pero sí lo son \emph{asintóticamente}: lejos del
origen, $\J_{t,s}\approx \J_{t-1,s-1}$. La figura
\ref{fig:estructura} lo muestra. Esa propiedad es la semilla de todo el
algoritmo rápido.

\section{El teorema de la función implícita}

Es la pieza que convierte ``resolver el modelo'' en ``invertir una matriz''.

Supongamos que el equilibrio está definido por un sistema
\begin{equation}
\mathbf{H}(\mathbf{U}, \mathbf{Z}) = \mathbf{0},
\end{equation}
donde $\mathbf{U}$ junta las trayectorias de las variables endógenas que
buscamos (las \emph{incógnitas}) y $\mathbf{Z}$ las de los choques exógenos. Si
$\mathbf{H}$ es diferenciable y su jacobiano respecto de $\mathbf{U}$ es
invertible en el estado estacionario, entonces alrededor de ese punto existe una
función implícita $\mathbf{U}(\mathbf{Z})$ y su derivada es
\begin{equation}
d\mathbf{U} = -\,\mathbf{H}_U^{-1}\,\mathbf{H}_Z\; d\mathbf{Z}.
\label{eq:ift}
\end{equation}

Nada más. Toda la solución del equilibrio general es esa línea. El trabajo del
método consiste, enteramente, en construir $\mathbf{H}_U$ y $\mathbf{H}_Z$.

\begin{trampa}
La invertibilidad de $\mathbf{H}_U$ no es un tecnicismo: es la condición de
\emph{determinación local} del modelo. Si el sistema está mal planteado --- por
ejemplo, si uno impone la ley de Walras y además todas las condiciones de
vaciado, de modo que una ecuación es redundante --- entonces $\mathbf{H}_U$ es
singular y el código falla con un error de álgebra lineal que no dice nada útil.
Si está casi singular, el código no falla pero devuelve basura. La sección
\ref{sec:diagnostico} explica cómo diagnosticarlo.
\end{trampa}

\section{Equivalencia cierta y la representación de medias móviles}

Esta es la justificación formal de por qué calcular una trayectoria determinista
alcanza para describir un modelo estocástico.

Cuando se linealiza un modelo alrededor de su estado estacionario, la solución
tiene la propiedad de \emph{equivalencia cierta}: las reglas de decisión óptimas
no dependen de la varianza de los choques. Un agente que enfrenta un choque
aleatorio se comporta exactamente igual que uno que enfrenta el valor esperado
de ese choque con certeza.

La consecuencia es poderosa. Sea $d\mathbf{X}_s$ la respuesta impulso de las
variables endógenas a una innovación unitaria que ocurre $s$ períodos después.
Entonces el proceso estocástico que sigue la economía es
\begin{equation}
d\tilde{X}_t = \sum_{s=0}^{T-1} d X_s\, \varepsilon_{t-s},
\label{eq:ma}
\end{equation}
que es una representación $\mathrm{MA}(T-1)$. Es decir: \emph{la función
impulso-respuesta es el modelo}. Con ella se puede simular, calcular momentos
poblacionales exactos y evaluar la verosimilitud, sin volver a tocar la
distribución ni el problema individual. El capítulo \ref{cap:estimacion}
explota esto a fondo.

\section{Choques MIT}

Un \emph{choque MIT} es un choque de probabilidad cero: la economía está en su
estado estacionario, los agentes creen que seguirá así para siempre, y en la
fecha $0$ ocurre algo inesperado. A partir de ese instante los agentes tienen
previsión perfecta sobre toda la trayectoria futura.

Es un dispositivo lógicamente incoherente --- los agentes son sorprendidos por
algo que asignaban probabilidad cero --- pero es exactamente la derivada que
necesitamos. Y por equivalencia cierta, a primer orden esa derivada coincide con
la solución del modelo genuinamente estocástico. La incoherencia no cuesta nada
mientras uno se quede en primer orden.

\section{Discretización: Markov, grillas y loterías}

El último ingrediente es técnico pero decisivo para la precisión.

\paragraph{Procesos exógenos} Un proceso AR(1) en logaritmos,
$\log z_t = \rho \log z_{t-1} + \sigma\epsilon_t$, se aproxima por una cadena de
Markov finita. El método de Rouwenhorst (1995) domina al de Tauchen cuando
$\rho$ es alto --- que es el caso típico en calibraciones trimestrales, donde
$\rho \approx 0.95$ o más --- porque replica exactamente la persistencia y la
varianza incondicional del proceso continuo para cualquier número de estados.

\paragraph{Grillas} El estado endógeno se discretiza sobre una grilla que debe
ser \emph{densa donde la política tiene curvatura}. En modelos de mercados
incompletos, eso es cerca de la restricción de endeudamiento. Una grilla
doblemente exponencial es la elección estándar.

\paragraph{La lotería de Young} Este es el punto que más errores silenciosos
causa. La política $a'(e,a_-)$ devuelve un número real que casi nunca cae sobre
un punto de la grilla. ¿Qué hacer? La respuesta ingenua --- redondear al punto
más cercano --- introduce un sesgo de agregación de primer orden que contamina
todos los jacobianos. La respuesta correcta (Young, 2010) es repartir la masa
entre los dos puntos adyacentes con pesos que \emph{preserven exactamente la
media}:
\begin{equation}
\text{si } a_i \le a' < a_{i+1}, \quad
p = \frac{a_{i+1}-a'}{a_{i+1}-a_i}
\;\Longrightarrow\;
p\, a_i + (1-p)\, a_{i+1} = a'.
\end{equation}
Con esto, la matriz de transición $\Lambda$ es una matriz de Markov legítima y
el ahorro agregado calculado sobre la grilla coincide con el ahorro agregado
verdadero.

\begin{trampa}
Hay dos operadores que usan los mismos índices y pesos de la lotería, pero en
\textbf{orden inverso}. El paso hacia adelante de la distribución aplica primero
la lotería y después la matriz de Markov exógena; el paso de los vectores de
expectativa aplica primero la matriz exógena y después la lotería. Formalmente,
si $\Lambda' = \Pi' L'$ entonces $\Lambda = L\,\Pi$. Invertir ese orden produce
jacobianos que se \emph{ven} razonables --- tienen la forma correcta, los signos
correctos, las magnitudes aproximadamente correctas --- y están mal. Es el error
más caro de depurar de todo el método. La única defensa confiable es la
verificación contra el método directo del capítulo \ref{cap:precision}.
\end{trampa}
